\section{Direct sums and the joint isometry condition}

\begin{lemma}[Block decomposition of the direct-sum transfer map]%
    \label{lem:block_diagonal_transfer_sum_to_block}
    \lean{MPSTensor.blockDiagonal'_transferSum_toBlock}
    \leanok
    \uses{def:mixed_transfer_rect}
    Let $(B_k)_k$ be a family of tensors and let $\bigoplus_k B_k$ be their
    direct sum on the total bond space. For every bond matrix $X$ and every pair
    of blocks, the $(j,j')$ bond block of the transfer sum
    \begin{align}
        \sum_i \left(\bigoplus_k B_k^i\right)X
            \left(\bigoplus_k B_k^i\right)^\dagger
        \notag
    \end{align}
    equals $\E_{j,j'}$ applied to the $(j,j')$ bond block of $X$.
\end{lemma}
\begin{proof}\leanok
    A block-diagonal factor connects only matching blocks, so the $(j,j')$ block
    of the product is
    \begin{align}
        \left(\sum_i \left(\bigoplus_k B_k^i\right)X
            \left(\bigoplus_k B_k^i\right)^\dagger\right)_{j,j'}
        &= \sum_i B_j^i X_{j,j'}(B_{j'}^i)^\dagger
         = \E_{j,j'}(X_{j,j'}).
        \notag
    \end{align}
\end{proof}

\begin{lemma}[Pairwise idempotence of the block transfer sum]%
    \label{lem:block_transfer_sum_idempotent_of_pairwise_mixed}
    \lean{MPSTensor.blockTransferSum_idempotent_of_pairwise_mixedTransferMap₂}
    \leanok
    \uses{def:mixed_transfer_rect, lem:block_diagonal_transfer_sum_to_block}
    Let $(B_k)_k$ be a family of tensors. If every mixed transfer operator
    $\E_{j,j'}$ is idempotent, then the block transfer sum $\mathcal S_B$ is
    idempotent: $\mathcal S_B(\mathcal S_B(X))=\mathcal S_B(X)$.
\end{lemma}
\begin{proof}\leanok
    By Lemma~\ref{lem:block_diagonal_transfer_sum_to_block}, the $(j,j')$
    block satisfies
    \begin{align}
        (\mathcal S_B^2(X))_{j,j'}
        &= \E_{j,j'}^2(X_{j,j'})
         = \E_{j,j'}(X_{j,j'})
         = (\mathcal S_B(X))_{j,j'}.
        \notag
    \end{align}
    Equality on every block gives the asserted equality on the total bond
    space.
\end{proof}

\begin{theorem}[Pairwise mixed-transfer criterion for a direct sum]%
    \label{thm:direct_sum_transfer_idempotent_iff_pairwise_mixed}
    \lean{MPSTensor.isTransferIdempotent_directSumTensor_iff_pairwise_mixedTransferMap₂_isIdempotentElem}
    \leanok
    \uses{def:mps_transfer_idempotence, def:mixed_transfer_rect,
        lem:block_diagonal_transfer_sum_to_block,
        lem:block_transfer_sum_idempotent_of_pairwise_mixed}
    Suppose every bond dimension is positive. The transfer map of
    $\bigoplus_k B_k$ is idempotent if and only if, for every pair $j,j'$,
    the mixed transfer operator is idempotent:
    \begin{align}
        \E_{\oplus_k B_k}^2=\E_{\oplus_k B_k}
        \quad\Longleftrightarrow\quad
        \E_{j,j'}^2=\E_{j,j'}.
        \notag
    \end{align}
\end{theorem}
\begin{proof}\leanok
    By Lemma~\ref{lem:block_diagonal_transfer_sum_to_block}, the $(j,j')$
    block of the direct-sum transfer map depends only on the $(j,j')$ block of
    its argument and is obtained by applying $\E_{j,j'}$. Thus pairwise
    idempotence gives
    \begin{align}
        (\E_{\oplus_k B_k}^2(X))_{j,j'}
        &= \E_{j,j'}^2(X_{j,j'})
         = \E_{j,j'}(X_{j,j'})
         = (\E_{\oplus_k B_k}(X))_{j,j'}
        \notag
    \end{align}
    for every pair of blocks, and hence gives idempotence on the total bond
    space, as in
    Lemma~\ref{lem:block_transfer_sum_idempotent_of_pairwise_mixed}.
    Conversely, every matrix $M$ of the appropriate rectangular size occurs as
    the $(j,j')$ block of some $Y$ on the total bond space. Applying
    Lemma~\ref{lem:block_diagonal_transfer_sum_to_block} and whole-tensor
    idempotence gives
    \begin{align}
        \E_{j,j'}^2(M)
        &= (\E_{\oplus_k B_k}^2(Y))_{j,j'}
         = (\E_{\oplus_k B_k}(Y))_{j,j'}
         = \E_{j,j'}(M).
        \notag
    \end{align}
\end{proof}

\begin{lemma}[Scaling of a rectangular mixed transfer operator]%
    \label{lem:mixed_transfer_rect_smul}
    \lean{MPSTensor.mixedTransferMap₂_smul}
    \leanok
    \uses{def:mixed_transfer_rect}
    For complex scalars $c,e$ and tensors $A,B$ of possibly different bond
    dimensions, $\E_{cA,eB}=c\overline e\,\E_{A,B}$.
\end{lemma}
\begin{proof}\leanok
    Conjugate transposition changes the second scalar to its complex conjugate:
    $\sum_i(cA^i)X(eB^i)^\dagger
    =c\overline e\sum_iA^iX(B^i)^\dagger$.
\end{proof}

\begin{theorem}[Phase coherence for literal repeated blocks]%
    \label{thm:literal_repeated_blocks_phase_coherence}
    \lean{MPSTensor.phases_eq_of_isTransferIdempotent_directSum_scaled_self}
    \leanok
    \uses{def:mps_transfer_idempotence, lem:mixed_transfer_rect_smul,
        thm:direct_sum_transfer_idempotent_iff_pairwise_mixed}
    Let $A$ have positive bond dimension and a nonzero idempotent transfer map.
    Let $\mu_q$ have unit modulus. If the literal bond direct sum
    $\bigoplus_q \mu_qA$ has an idempotent transfer map, then
    $\mu_q=\mu_{q'}$ for every pair $q,q'$.
\end{theorem}
\begin{proof}\leanok
    The pairwise criterion
    (Theorem~\ref{thm:direct_sum_transfer_idempotent_iff_pairwise_mixed})
    gives that $\E_{\mu_qA,\mu_{q'}A}$ is idempotent.
    Lemma~\ref{lem:mixed_transfer_rect_smul} and
    $\E_{A,A}=\E_A$ give
    $\E_{\mu_qA,\mu_{q'}A}
    =\mu_q\overline{\mu_{q'}}\E_{A,A}
    =\mu_q\overline{\mu_{q'}}\E_A$.
    Hence $(\mu_q\overline{\mu_{q'}})\E_A$ is idempotent. Since $\E_A$ is
    nonzero and
    $|\mu_q\overline{\mu_{q'}}|=|\mu_q||\mu_{q'}|=1$,
    \begin{align}
        (\mu_q\overline{\mu_{q'}})^2\E_A
        =\mu_q\overline{\mu_{q'}}\E_A
        &\quad\Longrightarrow\quad
        \mu_q\overline{\mu_{q'}}
        (\mu_q\overline{\mu_{q'}}-1)=0
        \quad\Longrightarrow\quad
        \mu_q\overline{\mu_{q'}}=1.
        \notag
    \end{align}
    Since $|\mu_{q'}|=1$ gives
    $\overline{\mu_{q'}}\mu_{q'}=1$, one has
    $\mu_q
    =\mu_q(\overline{\mu_{q'}}\mu_{q'})
    =(\mu_q\overline{\mu_{q'}})\mu_{q'}
    =\mu_{q'}$.

% Scope note: this is the corrected literal-block conclusion from
% \cite[Theorem~3.11, lines~543--563]{Cirac2016MPDO_arXiv}; the printed
% independent-phase statement requires a physical-space construction not
% specified by an equation in that passage.
\end{proof}

\begin{theorem}[Cross-block transfer vanishing for a direct-sum fixed point]%
    \label{thm:bnt_locally_orthogonal_of_rfp_direct_sum}
    \lean{MPSTensor.isBNTLocallyOrthogonal_of_isTransferIdempotent_directSum}
    \leanok
    \uses{def:is_bnt_locally_orthogonal, def:mps_transfer_idempotence,
        def:mixed_transfer_rect}
    Let $(A_j)_j$ be a family of distinct irreducible left-canonical blocks, no
    two of which are gauge-phase equivalent, and suppose the direct sum
    $\bigoplus_j A_j$ is a renormalization fixed point. Then for every pair of
    distinct components the mixed transfer operator vanishes,
    $\E_{j,j'}=\sum_iA_j^i\otimes\overline{A_{j'}^i}=0$.
    This is the cross-block ($\delta_{j,j'}$) content of the isometry condition
    of \cite[Theorem~3.11 and Corollary~3.12]{Cirac2016MPDO_arXiv} at the level of
    the normal-tensor blocks; converting it to the isometries $U_j$ is a
    further step left open.
% The U-level conversion is recorded in docs/paper-gaps/cpsv16_rfp_isometry_scope.tex.
\end{theorem}
\begin{proof}\leanok
    \uses{lem:block_diagonal_transfer_sum_to_block,
        thm:transfer_operator_gap_irr_tp,
        thm:transfer_operator_gap_rect_irr_tp,
        lem:idempotent_spectral_radius_zero}
    By Lemma~\ref{lem:block_diagonal_transfer_sum_to_block} the $(j,j')$ bond
    block of the direct-sum transfer map acts as $\E_{j,j'}$. Whole-tensor
    idempotence $\E^2=\E$ therefore restricts to each block as
    $\E_{j,j'}^2=\E_{j,j'}$. For distinct irreducible left-canonical blocks,
    Theorem~\ref{thm:transfer_operator_gap_irr_tp} gives
    $\rho_{\spec}(\E_{j,j'})<1$ when the bond dimensions agree, because the
    blocks are not gauge-phase equivalent;
    Theorem~\ref{thm:transfer_operator_gap_rect_irr_tp} gives the same
    inequality when the bond dimensions differ. Consequently,
    Lemma~\ref{lem:idempotent_spectral_radius_zero} gives
    \begin{align}
        \E_{j,j'}^2=\E_{j,j'},\qquad
        \rho_{\spec}(\E_{j,j'})<1
        \quad\Longrightarrow\quad
        \E_{j,j'}=0.
        \notag
    \end{align}
\end{proof}

\begin{definition}[Joint isometry condition]%
    \label{def:is_residual_isometry_family}
    \lean{MPSTensor.IsResidualIsometryFamily}
    \leanok
    A family $(U_j)_j$ satisfies the \emph{joint isometry condition} when
    \begin{align}
        \sum_i(U_j^i)_{\alpha,\beta}
            \overline{(U_{j'}^i)_{\alpha',\beta'}}
        &=\delta_{j,j'}\delta_{\alpha,\alpha'}\delta_{\beta,\beta'}.
        \notag
    \end{align}
    Each block satisfies the within-block pair-index orthonormality, and the
    cross-block sums between distinct blocks vanish. This is the isometry
    condition of \cite[Theorem~3.11]{Cirac2016MPDO_arXiv}, split into its $j=j'$
    and $j\ne j'$ cases.
\end{definition}

\begin{theorem}[One-letter Gram matrix under the joint isometry condition]%
    \label{thm:residual_isometry_family_one_gram}
    \lean{MPSTensor.IsResidualIsometryFamily.wordEntryFamily_one_gram}
    \leanok
    \uses{def:is_residual_isometry_family}
    Let $\mathcal V=\bigsqcup_j\{j\}\times[D_j]\times[D_j]$ be the disjoint
    union of the within-sector virtual pairs. For $x=(j,\alpha,\beta)$ in
    $\mathcal V$, set $u_x(i)=(U_j^i)_{\alpha,\beta}$. Then
    $\sum_i u_x(i)\overline{u_y(i)}=\delta_{x,y}$.
\end{theorem}

\begin{proof}\leanok
    If $x$ and $y$ belong to the same sector, the assertion is the
    within-sector part of the joint isometry condition. If they belong to
    distinct sectors, it is the cross-sector vanishing condition.
\end{proof}

\begin{lemma}[Covariance of the mixed transfer operator]%
    \label{lem:mixed_transfer_conj_apply}
    \lean{MPSTensor.mixedTransferMap₂_conj_apply}
    \leanok
    \uses{def:mixed_transfer_rect}
    For decompositions $A^i=X_AD_AU^iX_A^{-1}$ and
    $B^i=X_BD_BV^iX_B^{-1}$, the mixed transfer operator of $A,B$ is the
    mixed transfer operator of the tensors $U,V$ conjugated by the outer
    factors:
    \begin{align}
        \E_{A,B}(Y)
        &=X_AD_A\,\E_{U,V}\!\left(X_A^{-1}Y(X_B^{-1})^\dagger\right)(X_BD_B)^\dagger.
        \notag
    \end{align}
\end{lemma}
\begin{proof}\leanok
    Substitute the two decompositions into
    $\E_{A,B}(Y)=\sum_iA^iY(B^i)^\dagger$ and collect the outer factors,
    which are independent of the summation index, outside the sum.
\end{proof}

\begin{lemma}[Vanishing after removing the diagonal factors]%
    \label{lem:mixed_transfer_eq_zero_of_conj}
    \lean{MPSTensor.mixedTransferMap₂_eq_zero_of_conj}
    \leanok
    \uses{def:mixed_transfer_rect}
    With decompositions $A^i=X_AD_AU^iX_A^{-1}$ and
    $B^i=X_BD_BV^iX_B^{-1}$ whose factors $X_A,D_A,X_B,D_B$ are
    invertible, if $\E_{A,B}=0$ then $\E_{U,V}=0$.
\end{lemma}
\begin{proof}\leanok
    \uses{lem:mixed_transfer_conj_apply}
    The covariance identity Lemma~\ref{lem:mixed_transfer_conj_apply} expresses
    $\E_{A,B}$ as $\E_{U,V}$ conjugated by the invertible outer factors
    $X_AD_A$ and $(X_BD_B)^\dagger$, so $\E_{A,B}=0$ forces $\E_{U,V}=0$.
\end{proof}

\begin{lemma}[Entrywise cross-block isometry sum]%
    \label{lem:residual_isometry_entry}
    \lean{MPSTensor.residual_isometry_entry_of_mixedTransferMap₂_eq_zero}
    \leanok
    \uses{def:mixed_transfer_rect}
    If $\E_{U,V}=0$, then, for all virtual indices,
    \begin{align}
        \sum_i(U^i)_{\alpha,\beta}
        \overline{(V^i)_{\alpha',\beta'}}
        &=0.
        \notag
    \end{align}
\end{lemma}
\begin{proof}\leanok
    Apply $\E_{U,V}=0$ to the matrix unit supported at $(\beta,\beta')$ and read
    the $(\alpha,\alpha')$ entry.
\end{proof}

\begin{lemma}[Each block of a direct-sum fixed point is a fixed point]%
    \label{lem:isrfp_block_of_rfp_direct_sum}
    \lean{MPSTensor.isTransferIdempotent_block_of_isTransferIdempotent_directSum}
    \leanok
    \uses{def:mps_transfer_idempotence, def:mixed_transfer_rect}
    If the direct sum $\bigoplus_k B_k$ is a renormalization fixed point, with
    $\dim_k\ge1$ for all $k$, then each block $B_j$ is a renormalization fixed
    point.
\end{lemma}
\begin{proof}\leanok
    \uses{lem:block_diagonal_transfer_sum_to_block}
    The diagonal mixed transfer operator $\E_{j,j}$ is the transfer map of
    $B_j$, and whole-tensor idempotence makes it idempotent, which is exactly the
    renormalization fixed-point condition for $B_j$.
\end{proof}

\begin{theorem}[Joint isometry condition for a direct-sum fixed point]%
    \label{thm:residual_isometry_of_rfp_direct_sum}
    \lean{MPSTensor.exists_residualIsometryFamily_of_isTransferIdempotent_directSum}
    \leanok
    \uses{def:is_residual_isometry_family, def:mps_transfer_idempotence,
        def:is_isometry_canonical_form}
    Let $(A_j)_j$ be a family of normal, irreducible, left-canonical blocks, no
    two of which are gauge-phase equivalent, with $\dim_j\ge1$ for all $j$, and
    suppose the direct sum $\bigoplus_j A_j$ is a renormalization fixed point.
    Then there are invertible matrices $X_j$, trace-one positive diagonal
    matrices $\rho_j=\diag(\lambda_{j,\alpha})$, and tensors $U_j$ with
    \begin{align}
        A_j^i&=X_j\sqrt{\rho_j}\,U_j^iX_j^{-1},
        \notag
    \end{align}
    such that $(U_j)_j$ satisfies the joint isometry condition. These are the
    isometry equations of \cite[Corollary~3.12]{Cirac2016MPDO_arXiv} for the
    direct sum of distinct normal-tensor blocks.
% Scope restriction: the Lean theorem takes the per-block normality,
% irreducibility, left-canonical condition, and gauge-phase distinctness as
% explicit hypotheses; extracting these from a single canonical-form RFP
% predicate is the remaining step
% (docs/paper-gaps/cpsv16_rfp_isometry_scope.tex).
\end{theorem}
\begin{proof}\leanok
    \uses{thm:bnt_locally_orthogonal_of_rfp_direct_sum,
        thm:isometry_canonical_form_of_rfp_nt, lem:mixed_transfer_conj_apply,
        lem:isrfp_block_of_rfp_direct_sum, lem:mixed_transfer_eq_zero_of_conj,
        lem:residual_isometry_entry}
    Whole-tensor idempotence makes the diagonal mixed transfer operator of each
    block its own transfer map, hence each block is a renormalization fixed
    point; the square-root fixed-point form
    (Theorem~\ref{thm:isometry_canonical_form_of_rfp_nt}) supplies the
    decomposition $A_j^i=X_j\sqrt{\rho_j}U_j^iX_j^{-1}$ with the within-block
    orthonormality, which is the $j=j'$ case. For $j\ne j'$,
    Theorem~\ref{thm:bnt_locally_orthogonal_of_rfp_direct_sum} gives
    $\E_{A_j,A_{j'}}=0$; by the covariance
    Lemma~\ref{lem:mixed_transfer_conj_apply} and invertibility of the outer
    factors $\E_{U_j,U_{j'}}=0$, and reading the entry of this operator at a
    matrix unit yields the cross-block sum
    $\sum_i(U_j^i)_{\alpha,\beta}
    \overline{(U_{j'}^i)_{\alpha',\beta'}}=0$.
\end{proof}

\begin{theorem}[Direct sum of square-root fixed-point blocks]%
    \label{thm:is_rfp_of_isometry_canonical_form_direct_sum}
    \lean{MPSTensor.isTransferIdempotent_directSumTensor_of_isIsometryCanonicalForm}
    \leanok
    \uses{def:is_isometry_canonical_form, def:mps_transfer_idempotence,
        def:mixed_transfer_rect}
    Let $(B_k)_k$ be a family of tensors, each in the square-root fixed-point
    form, whose cross-block mixed transfer operators vanish,
    $\E_{j,j'}=0$ for $j\ne j'$. Then
    the direct sum $\bigoplus_k B_k$ is a renormalization fixed point. This is the
    backward direction of the structural characterization of pure-state
    renormalization fixed points \cite[Section~3]{Cirac2016MPDO_arXiv}, in the
    distinct-blocks case where each normal tensor appears once; the cross-block
    vanishing is the $j\ne j'$ part of the isometry condition of that
    characterization, so it is part of the source joint isometry condition.
% Scope restriction (multiplicity-one): the source canonical form allows each
% normal tensor to repeat with a multiplicity and a phase; this is the
% single-copy case (docs/paper-gaps/cpsv16_rfp_isometry_scope.tex).
\end{theorem}
\begin{proof}\leanok
    \uses{thm:is_rfp_of_isometry_canonical_form,
        lem:block_diagonal_transfer_sum_to_block}
    The direct-sum transfer map decouples block by block: by
    Lemma~\ref{lem:block_diagonal_transfer_sum_to_block} its $(j,j')$ bond block
    acts as $\E_{j,j'}$ on the $(j,j')$ bond block of the argument. Each diagonal
    block $\E_{j,j}$ is the transfer map of $B_j$, which is idempotent because
    $B_j$ is a renormalization fixed point
    (Theorem~\ref{thm:is_rfp_of_isometry_canonical_form}); each off-diagonal block
    $\E_{j,j'}$ with $j\ne j'$ vanishes by hypothesis. Therefore, for every bond
    block,
    \begin{align}
        [\E_{\bigoplus B}^2(Y)]_{j,j'}
        &=\E_{j,j'}\!(\E_{j,j'}(Y_{j,j'}))
         =\E_{j,j'}(Y_{j,j'})
         =[\E_{\bigoplus B}(Y)]_{j,j'},
        \notag
    \end{align}
    so $\E_{\bigoplus B}^2=\E_{\bigoplus B}$ and the direct
    sum is a renormalization fixed point.
\end{proof}

\begin{theorem}[Distinct-block fixed-point characterization]%
    \label{thm:rfp_directSumTensor_iff}
    \lean{MPSTensor.isTransferIdempotent_directSumTensor_iff}
    \leanok
    \uses{def:mps_transfer_idempotence, def:is_isometry_canonical_form,
        def:mixed_transfer_rect,
        def:normal, def:irreducible_tensor, def:gauge_phase_equiv}
    Let $(B_k)_k$ be a family of normal, irreducible, left-canonical blocks with
    $\dim_k\ge1$ for all $k$, no two of which are gauge-phase equivalent. Then
    the direct sum $\bigoplus_k B_k$ is a renormalization fixed point if and only
    if each block is in the square-root fixed-point form and the mixed transfer
    operators between distinct blocks vanish:
    \begin{align}
        \left(\bigoplus_k B_k\text{ is a RFP}\right)
        &\iff
        \left(\forall k,\ B_k\text{ in square-root form}\right)
        \wedge
        \left(\forall j\ne j',\ \E_{j,j'}=0\right).
        \notag
    \end{align}
    This is the distinct-blocks (multiplicity-one, phase-one) case of the
    structural characterization of pure-state renormalization fixed points
    \cite[Theorem~3.11]{Cirac2016MPDO_arXiv}, combining its forward and backward
    directions.
% Scope restriction (multiplicity-one): the source canonical form
% A^i = \oplus_j \oplus_q mu_{j,q} X_{j,q} Lambda_j U^i_j X_{j,q}^{-1} allows
% each normal tensor to repeat with multiplicity r_j and a unit-modulus phase
% mu_{j,q}; this is the single-copy case
% (docs/paper-gaps/cpsv16_rfp_isometry_scope.tex).
\end{theorem}
\begin{proof}\leanok
    \uses{thm:isometry_canonical_form_of_rfp_nt,
        lem:isrfp_block_of_rfp_direct_sum,
        thm:bnt_locally_orthogonal_of_rfp_direct_sum,
        thm:is_rfp_of_isometry_canonical_form_direct_sum}
    For the forward direction, whole-tensor idempotence makes each diagonal mixed
    transfer operator the transfer map of its block
    (Lemma~\ref{lem:isrfp_block_of_rfp_direct_sum}), so each block is a
    renormalization fixed point and hence in the square-root fixed-point form
    (Theorem~\ref{thm:isometry_canonical_form_of_rfp_nt}), while the off-diagonal
    operators vanish, $\E_{j,j'}=0$ for $j\ne j'$, by
    Theorem~\ref{thm:bnt_locally_orthogonal_of_rfp_direct_sum}. Conversely, a
    direct sum of square-root fixed-point blocks whose cross-block operators
    vanish is a renormalization fixed point
    (Theorem~\ref{thm:is_rfp_of_isometry_canonical_form_direct_sum}).
\end{proof}

\begin{corollary}[BNT basis direct-sum characterization]%
    \label{cor:rfp_bnt_basis_direct_sum_iff}
    \lean{MPSTensor.IsBNTCanonicalForm.isTransferIdempotent_basisDirectSum_iff}
    \lean{MPSTensor.IsBNTCanonicalForm.exists_residualIsometryFamily_of_isTransferIdempotent_basisDirectSum}
    \leanok
    \uses{def:sector_bnt_cf, def:mps_transfer_idempotence,
        def:is_isometry_canonical_form, def:mixed_transfer_rect,
        def:is_residual_isometry_family}
    Let $P$ satisfy the auxiliary BNT sector hypotheses, with distinct basis tensors $(B_j)_j$.
    Then the direct sum containing one copy of each basis tensor is a
    renormalization fixed point if and only if every $B_j$ has the square-root
    fixed-point form and $\E_{j,j'}=0$ for $j\ne j'$.
    This is the multiplicity-one, phase-one specialization of
    \cite[Theorem~3.11 and Corollary~3.12]{Cirac2016MPDO_arXiv}. It concerns
    the direct sum of the basis representatives, not the repeated-copy tensor
    carrying the weights $\mu_{j,q}$.

    If the basis direct sum is a renormalization fixed point, there are also
    invertible matrices $X_j$, trace-one positive diagonal matrices
    $\rho_j=\diag(\lambda_{j,\alpha})$, and tensors $U_j$ such that
    $B_j^i=X_j\sqrt{\rho_j}\,U_j^iX_j^{-1}$, and the family $(U_j)_j$
    satisfies the joint isometry condition, including the cross-block
    equations for $j\ne j'$.
% The physical-space construction of the repeated-copy index is recorded in
% docs/paper-gaps/cpsv16_rfp_isometry_scope.tex.
\end{corollary}
\begin{proof}\leanok
    \uses{thm:rfp_directSumTensor_iff,
        thm:postblocked_representative_grouped_insert_eq,
        thm:primitive_iff_period_one,
        thm:normal_tp_primitive_irred}
    Positive basis dimensions supply the nonzero bond spaces. For each basis
    tensor, irreducibility and left-canonical normalization are part of the
    auxiliary BNT sector hypotheses. If $m_j>0$ is the peripheral period of
    $B_j$, the
    normalized self-overlap hypothesis and the periodic self-overlap formula
    used in Theorem~\ref{thm:postblocked_representative_grouped_insert_eq} give,
    along the same subsequence,
    \begin{align}
        \lim_{k\to\infty}O_{B_jB_j}(m_jk)&=1,
        \notag\\
        \lim_{k\to\infty}O_{B_jB_j}(m_jk)&=m_j.
        \notag
    \end{align}
    Hence $m_j=1$. The transfer map is primitive by
    Theorem~\ref{thm:primitive_iff_period_one}, and normality follows from
    Theorem~\ref{thm:normal_tp_primitive_irred}. Gauge-phase distinctness is
    also part of the auxiliary BNT sector hypotheses. The equivalence now
    follows from
    Theorem~\ref{thm:rfp_directSumTensor_iff}; the same derived hypotheses in
    Theorem~\ref{thm:residual_isometry_of_rfp_direct_sum} gives the joint
    isometry condition.
\end{proof}


\begin{lemma}[Transfer idempotence under coisometric reconstruction]%
    \label{lem:rfp_coisometry_reconstruction_iff}
    \lean{MPSTensor.isTransferIdempotent_coisometry_reconstruction_iff}
    \leanok
    \uses{def:mps_transfer_idempotence}
    Let $A$ and $B$ be tensors with bond dimensions $D$ and $R$. Suppose that
    $U\in\C^{R\times D}$ satisfies $UU^\dagger=I_R$ and
    \begin{align}
        A^i&=U^\dagger B^iU
        \qquad\text{for every }i.
        \notag
    \end{align}
    Then $A$ is a renormalization fixed point if and only if $B$ is a
    renormalization fixed point.
\end{lemma}
\begin{proof}\leanok
    For every bond matrix $X$,
    \begin{align}
        \E_A(X)
        &=U^\dagger\E_B(UXU^\dagger)U.
        \notag
    \end{align}
    The identity $UU^\dagger=I_R$ gives
    $U(U^\dagger YU)U^\dagger=Y$. Applying these two identities twice proves
    each implication.
\end{proof}

\begin{lemma}[A fixed nonzero scaling of a normal tensor]%
    \label{lem:normal_scaled_rfp}
    \lean{MPSTensor.norm_eq_one_and_isTransferIdempotent_of_isNormalTensor_smul}
    \leanok
    \uses{def:nt_cpsv, def:mps_transfer_idempotence}
    Let $B$ be a normal tensor and let $c\in\C\setminus\{0\}$. If $cB$ is a
    renormalization fixed point, then
    \begin{align}
        |c|&=1,
        &\E_B^2&=\E_B.
        \notag
    \end{align}
\end{lemma}
\begin{proof}\leanok
    Normality gives $r(\E_B)=1$, while
    $\E_{cB}=|c|^2\E_B$. Since $c\ne0$, this operator is nonzero. A nonzero
    idempotent has spectral radius one, and hence
    \begin{align}
        1=r(\E_{cB})=|c|^2r(\E_B)=|c|^2.
        \notag
    \end{align}
    Thus $|c|=1$, so $\E_{cB}=\E_B$ and the idempotence of $\E_B$ follows.
\end{proof}

\begin{theorem}[Active blocks of a CPSV fixed point]%
    \label{thm:cpsv_active_weight_norm_one_and_block_rfp}
    \lean{MPSTensor.CPSVCanonicalFormData.active_weight_norm_one_and_block_rfp}
    \leanok
    \uses{def:cpsv_canonical_form, def:mps_transfer_idempotence,
        lem:rfp_coisometry_reconstruction_iff, lem:normal_scaled_rfp,
        lem:isrfp_block_of_rfp_direct_sum}
    Suppose $A$ has CPSV canonical-form data
    \begin{align}
        A^i
        &=U^\dagger\!\left(\bigoplus_{k=1}^r\mu_kA_k^i\right)U,
        &UU^\dagger&=I,
        \notag
    \end{align}
    where every $A_k$ is normal. If $A$ is a renormalization fixed point, then
    every active block $k$, characterized by $\mu_k\ne0$, satisfies
    \begin{align}
        |\mu_k|&=1,
        &\E_{A_k}^2&=\E_{A_k}.
        \notag
    \end{align}
\end{theorem}
\begin{proof}\leanok
    \uses{lem:rfp_coisometry_reconstruction_iff,
        lem:isrfp_block_of_rfp_direct_sum, lem:normal_scaled_rfp}
    Lemma~\ref{lem:rfp_coisometry_reconstruction_iff} makes the retained direct
    sum $\bigoplus_k\mu_kA_k$ a renormalization fixed point.
    Lemma~\ref{lem:isrfp_block_of_rfp_direct_sum} then makes each scaled block
    $\mu_kA_k$ a renormalization fixed point. For an active block,
    Lemma~\ref{lem:normal_scaled_rfp} gives both conclusions.
\end{proof}

\begin{theorem}[An arbitrary BNT need not possess a joint residual isometry]%
    \label{thm:cpsv_cor312_arbitrary_bnt_counterexample}
    \lean{MPSTensor.cpsvCorollary312_arbitraryBNT_counterexample}
    \leanok
    \uses{def:cpsv_canonical_form, def:bnt_cpsv,
        def:mps_transfer_idempotence, def:is_residual_isometry_family}
    There is a canonical-form renormalization fixed point $A$ and a basis of
    normal tensors $(A_0,A_1)$ for it such that both $A_0$ and $A_1$ are
    renormalization fixed points, but the family admits no trace-normalized
    square-root decomposition with a joint residual isometry. One may take
    bond dimension one and
    \begin{align}
        A^0=A_0^0&=1,
        &A^1=A_0^1&=0,
        &A_1^0=A_1^1&=\frac{1}{\sqrt2}.
        \notag
    \end{align}
    The BNT coefficients are $1$ on $A_0$ and $0$ on $A_1$ at every length.
\end{theorem}
\begin{proof}\leanok
    The two matrix product vectors are linearly independent at every positive
    length, so $(A_0,A_1)$ is a BNT for $A$ with the displayed coefficients.
    Both transfer maps are the identity. In bond dimension one, trace
    normalization fixes the positive diagonal factor to $1$, and scalar
    conjugation has no effect. Any residual tensors would therefore equal
    $A_0$ and $A_1$. Their cross inner product is
    \begin{align}
        \sum_{i=0}^1 A_0^i\overline{A_1^i}
        &=\frac{1}{\sqrt2}\ne0,
        \notag
    \end{align}
    contradicting the off-diagonal joint-isometry equation.
\end{proof}



% CPSV16 Theorem 3.11 (`thm:charact-MPS`), lines 543--562.
% No Lean theorem currently derives this repeated-copy statement from literal
% `IsCPSVCanonicalForm`; the proved direct-sum result below the source level has
% one unit-weight copy of each distinct block.  The missing q-routing data are
% the documented source obstruction in closed issue #2598 and
% docs/paper-gaps/cpsv16_rfp_isometry_scope.tex.
% Formalization note: `CPSVCanonicalFormData` permits an omitted zero ambient
% bond complement, but this implementation detail is not added to the source
% theorem as a new hypothesis.
\begin{theorem}[Square-root characterization of fixed points]
    \label{thm:cpsv_charact_mps_status}
    \notready
    \uses{def:cpsv_canonical_form, def:bnt_cpsv,
        def:mps_transfer_idempotence, def:is_isometry_canonical_form,
        def:is_residual_isometry_family, thm:phase_flip_tensor_not_rfp}
    CPSV asserts that a tensor $A$ in canonical form is a renormalization fixed
    point if and only if it can be written as
    \begin{align}
        A^i
        &=\bigoplus_{j=1}^g\bigoplus_{q=1}^{r_j}
          \mu_{j,q}X_{j,q}\sqrt{\rho_j}\,U_j^iX_{j,q}^{-1},
        \notag
    \end{align}
    where $|\mu_{j,q}|=1$, each $\rho_j$ is positive diagonal with
    $\tr(\rho_j)=1$, and
    \begin{align}
        \sum_i(U_j^i)_{\alpha,\beta}
          \overline{(U_{j'}^i)_{\alpha',\beta'}}
        &=\delta_{j,j'}\delta_{\alpha,\alpha'}\delta_{\beta,\beta'}.
        \notag
    \end{align}
    The source specifies that $\bigoplus_{j,q}$ is a simultaneous direct sum in
    the physical and virtual spaces, not an ordinary virtual block diagonal at
    each fixed physical letter $i$ \cite[Theorem~3.11,
    lines~561--562]{Cirac2016MPDO_arXiv}. It does not provide a coordinate map
    for the physical routing of $q$, a corresponding dimension hypothesis, or
    a $q$-indexed isometry equation. Consequently the display does not yet
    determine a source-faithful coordinate-level predicate. Under the naive
    fixed-$i$ virtual-block interpretation, the claimed converse is false by
    Theorem~\ref{thm:phase_flip_tensor_not_rfp}. The routed characterization
    therefore remains not ready pending a source clarification or erratum.

    \textbf{Local fix (square-root diagonal).} CPSV prints a bare $\Lambda_j$
    in Theorem~3.11, whereas the reference tensor at line~1300 has coefficients
    $\sqrt{(\rho_j)_{\alpha,\alpha}}$, which forces the factor $\sqrt{\rho_j}$
    used above. This repair is documented in
    \path{docs/paper-gaps/cpsv16_rfp_isometry_scope.tex}.
\end{theorem}

% CPSV16 Corollary 3.12 (`III_cor3`), lines 583--590.
% The arbitrary-BNT statement is refuted by an inactive BNT member.  The active
% diagonal conclusion is proved for literal CPSV canonical-form data.  The
% joint residual family for active phase-class representatives is still open.
\begin{corollary}[Basis tensors of a fixed point]
    \label{cor:cpsv_iii_cor3_status}
    \notready
    \uses{def:cpsv_canonical_form, def:bnt_cpsv,
        def:mps_transfer_idempotence, def:is_isometry_canonical_form,
        def:is_residual_isometry_family,
        thm:cpsv_active_weight_norm_one_and_block_rfp,
        thm:cpsv_cor312_arbitrary_bnt_counterexample,
        thm:residual_isometry_of_rfp_direct_sum}
    Let $A$ be a CPSV canonical-form renormalization fixed point, and let
    $(A_j)_j$ be an arbitrary basis of normal tensors for $A$. CPSV asserts that
    there are invertible matrices $X_j$, positive diagonal matrices $\rho_j$
    with $\tr(\rho_j)=1$, and tensors $U_j$ such that
    \begin{align}
        A_j^i
        &=X_j\sqrt{\rho_j}\,U_j^iX_j^{-1},
        \notag\\
        \sum_i(U_j^i)_{\alpha,\beta}
          \overline{(U_{j'}^i)_{\alpha',\beta'}}
        &=\delta_{j,j'}\delta_{\alpha,\alpha'}\delta_{\beta,\beta'}.
        \notag
    \end{align}
    The assertion is false for an arbitrary BNT. Such a BNT may contain a
    normal tensor whose coefficient vanishes at every positive length, and
    Theorem~\ref{thm:cpsv_cor312_arbitrary_bnt_counterexample} shows that an
    unused member can violate the cross-block equation.

    The diagonal conclusion has a corrected active form. For the normal blocks
    $(A_k)_k$ occurring with nonzero weights $\mu_k$ in literal CPSV
    canonical-form data, Theorem~\ref{thm:cpsv_active_weight_norm_one_and_block_rfp}
    proves
    \begin{align}
        |\mu_k|&=1,
        &\E_{A_k}^2&=\E_{A_k}.
        \notag
    \end{align}
    The single-block fixed-point characterization then gives the
    trace-normalized square-root decomposition of each active block. It remains
    to prove the off-diagonal joint-isometry equations simultaneously for the
    representatives of the active gauge-phase classes. Theorem~\ref{thm:residual_isometry_of_rfp_direct_sum}
    proves these equations once the direct sum of those representatives is
    known to be a renormalization fixed point.

    \textbf{Local fix (square-root diagonal).} CPSV prints a bare $\Lambda_j$
    in Corollary~3.12, whereas the reference tensor at line~1300 has
    coefficients $\sqrt{(\rho_j)_{\alpha,\alpha}}$, which forces the factor
    $\sqrt{\rho_j}$ used above. The source defect and the corrected active
    statement are documented in
    \path{docs/paper-gaps/cpsv16_rfp_isometry_scope.tex}.
\end{corollary}

\begin{theorem}[Ambiguity of the literal repeated-phase display]%
    \label{thm:phase_flip_tensor_not_rfp}
    \lean{MPSTensor.phaseFlipTensor_literal_display_ambiguity}
    \leanok
    \uses{def:mps_transfer_idempotence, def:is_isometry_canonical_form}
    Let $B$ be the one-letter, one-dimensional tensor $B^0=(1)$, and form two
    block-diagonal copies with phases $1$ and $-1$:
    \begin{align}
        A^0&=\begin{pmatrix}1&0\\0&-1\end{pmatrix}.
        \notag
    \end{align}
    The representative $B$ is in the square-root fixed-point form, and both copy
    coefficients have unit modulus. Thus $A$ satisfies the literal virtual
    block-diagonal reading of the displayed repeated-copy formula in
    \cite[Section~3]{Cirac2016MPDO_arXiv}. It does not, however, satisfy the
    accompanying interpretation in which the copy index also belongs to the
    physical direct sum: there is no scalar $v$ such that
    $(A^0)^2=vA^0$. In particular, $A$ is not a renormalization fixed point
    in the sense of Definition~\ref{def:mps_transfer_idempotence}.
% Source ambiguity: equation III_CFI_RFP at lines 543--554 admits this literal
% virtual-block substitution, while lines 559--563 require j,q to contribute
% to both physical and virtual direct sums; see
% docs/paper-gaps/cpsv16_rfp_isometry_scope.tex.
\end{theorem}
\begin{proof}\leanok
    Taking the witness data of
    Definition~\ref{def:is_isometry_canonical_form} to be
    $X=\rho=U^0=(1)$, the scalar representative is in square-root fixed-point
    form, and $|1|=|-1|=1$.
    Direct multiplication gives $(A^0)^2=I$. Comparing diagonal entries in
    $I=vA^0$ would give simultaneously $v=1$ and $v=-1$, which is impossible.
    On the off-diagonal matrix unit $E_{01}$, the transfer map is conjugation by
    $A^0$, so $\E_A(E_{01})=-E_{01}$ and $\E_A^2(E_{01})=E_{01}$.
    Hence $\E_A^2\ne\E_A$.
\end{proof}
